\documentclass[twocolumn]{revtex4-2}

\usepackage{amsmath,amssymb,amsthm}
\usepackage{graphicx}
\usepackage{booktabs}
\usepackage{xcolor}
\usepackage{hyperref}

\newtheorem{theorem}{Theorem}

\newcommand{\Nf}{N_f}
\newcommand{\LambdaUV}{\Lambda_{\mathrm{UV}}}
\newcommand{\Lambdaemerg}{\Lambda^{\mathrm{emerg}}}
\newcommand{\Lambdaobs}{\Lambda_{\mathrm{obs}}}
\newcommand{\MPl}{M_{\mathrm{Pl}}}
\newcommand{\aADW}{\alpha_{\mathrm{ADW}}}
\newcommand{\GN}{G_{N}}
\newcommand{\GNem}{G_{N}^{\mathrm{emerg}}}
\newcommand{\GNsak}{G_{N}^{\mathrm{Sakharov}}}

% Auto-generated by update_counts.py — 2026-04-24T19:44:03.140400
% Do not edit manually. Run: uv run python scripts/update_counts.py
\newcommand{\totaltheorems}{3993}
\newcommand{\substantivetheorems}{3882}
\newcommand{\placeholdertheorems}{111}
\newcommand{\axiomcount}{1}
\newcommand{\sorrycount}{0}
\newcommand{\sorrypercent}{0.0}
\newcommand{\leanmodules}{167}
\newcommand{\leandefinitions}{3297}
\newcommand{\pythonmodules}{68}
\newcommand{\testfiles}{59}
\newcommand{\totaltests}{2836}
\newcommand{\figurecount}{111}
\newcommand{\notebookcount}{54}
\newcommand{\papercount}{19}
\newcommand{\aristotleproved}{322}
\newcommand{\aristotleruns}{44}
% --- Paper 18 / Phase 5t per-section counts (FermiHubbardDimer.lean) ---
\newcommand{\fhdTotal}{143}
\newcommand{\fhdWaveTwo}{13}
\newcommand{\fhdWaveThree}{6}
\newcommand{\fhdWaveFour}{22}
\newcommand{\fhdWaveFive}{22}
\newcommand{\fhdWaveSixABC}{38}
\newcommand{\fhdWaveSixExtra}{15}
\newcommand{\fhdWaveSeven}{19}
\newcommand{\fhdWaveEight}{8}


\begin{document}

\title{The cosmological constant in emergent form: a Lean-formalized
  Decision Gate on the microscopic-to-macroscopic match in the ADW
  heat-kernel programme}

\author{John G. Roehm}
\affiliation{Independent Researcher; \texttt{NetRxn137@gmail.com}}

\date{\today}

\begin{abstract}
We formalize, in the Lean~4 proof assistant, the
microscopic-to-macroscopic coefficient match for the ADW
emergent-gravity programme: closed-form expressions for the emergent
Newton constant $\GN$, emergent cosmological constant
$\Lambdaemerg$, and Stelle-basis higher-curvature triple
$(\alpha,\beta,\gamma)$ in microscopic parameters
$(\LambdaUV,\Nf)$. The substantive
\emph{Decision-Gate-E.4} result is a quantitative theorem: at the
natural microscopic point $(\LambdaUV \simeq \MPl,\ \Nf = 16)$, the
heat-kernel-induced
$\Lambdaemerg(\LambdaUV,\Nf) = a_{0}(\Nf)\cdot\LambdaUV^{4}$
exceeds the Planck-2018 observed value
$\Lambdaobs \simeq 2.6\times 10^{-47}\,\mathrm{GeV}^{4}$ by more than
$10^{120}$. The classical cosmological-constant problem is
\emph{reproduced in the emergent-gravity formulation}; no fine-tuning
of $(\LambdaUV,\Nf)$ resolves it within the natural high-energy
parameter band. The Lean proof is closed and zero-sorry; we ship a
companion Python pipeline reproducing the natural-parameter scan and
a verdict-band figure marking the resolved/intermediate/reproduced
phases. Together with Wave~1's Decision Gate~E.2 (heat-kernel $\leftrightarrow$
linearized-EFE Newton-constant calibration), this closes the
microscopic-coefficient match question: the ADW programme provides
\emph{no shelter} from the cosmological-constant fine-tuning problem.
\end{abstract}

\maketitle

\section{Introduction}

The ADW (Alexandrov-Diakonov-Vergeles) emergent-gravity programme
\cite{Sakharov1968,Adler1982,Diakonov2011,VladimirovDiakonov2012,Vergeles2025}
posits that the Einstein-Hilbert effective action arises from
integrating out a microscopic 8-fermion sector at a hard ultraviolet
cutoff $\LambdaUV$. At the trace level of the variational EFE
(companion paper, in preparation \cite{Roehm2026Wave4}), the
leading Seeley--DeWitt heat-kernel
coefficients $a_{0}, a_{2}$ produce, respectively, a cosmological-
constant scale $\Lambdaemerg \sim a_{0}(\Nf)\cdot\LambdaUV^{4}$ and
the Sakharov-Adler induced Newton constant
$\GNsak = 12\pi/(\Nf\,\LambdaUV^{2})$ \cite{Vassilevich2003}.

The natural question is whether the emergent
$\Lambdaemerg$ matches the Planck-2018 observed cosmological
constant $\Lambdaobs \simeq 2.6\times 10^{-47}\,\mathrm{GeV}^{4}$
\cite{Planck2018}. The classical cosmological-constant problem ---
that natural-cutoff field-theoretic estimates exceed $\Lambdaobs$ by
many orders of magnitude --- is well known; the question for an
emergent-gravity programme is whether the emergent reformulation
provides any \emph{shelter}. Two outcomes are publishable; both are
substantive:

\begin{enumerate}
\item \emph{Resolution}: under natural microscopic parameters,
$\Lambdaemerg$ approximately reproduces $\Lambdaobs$. This would
constitute a derivation of the cosmological constant from
microphysics --- a major theoretical result.
\item \emph{Reproduction}: $\Lambdaemerg$ exceeds $\Lambdaobs$ by
many orders of magnitude even in the natural-cutoff band,
reproducing the cosmological-constant problem in emergent form.
\end{enumerate}

We formalize both outcomes as a categorical \emph{Decision Gate~E.4}
verdict and prove the quantitative claim in the Lean~4 proof
assistant: the natural-parameter point lands strictly in the
``reproduction'' phase, with $\Lambdaemerg > 10^{100}\Lambdaobs$ at
$\LambdaUV \simeq 1.2\times 10^{19}\,\mathrm{GeV},\ \Nf = 16$.
The companion Python pipeline reproduces the corresponding ratio
(approximately $10^{122}$) numerically and labels the verdict band
across the full $(\LambdaUV,\Nf)$ plane.

\section{Microscopic predictions}

\subsection{Definitions}

Following Wave~1 \cite{Roehm2026Wave1}, the leading Seeley--DeWitt
coefficient for $\Nf$ free Dirac fermions in 4D vacuum is
\begin{equation}
  a_{0}(\Nf) = \frac{4\Nf}{(4\pi)^{2}},
  \label{eq:a0}
\end{equation}
matching Vassilevich \cite{Vassilevich2003} Eq.~(4.37) up to the
spinor-trace convention. Concretely, Eq.~(4.37) gives the
single-Dirac-fermion vacuum coefficient $a_{0}^{(\mathrm{single})}(x)
= \mathrm{tr}(\mathbb{1})/(4\pi)^{2} = 4/(4\pi)^{2}$ under the
convention $\mathrm{tr}(\mathbb{1}) = 4$ over the four-component
Dirac matrices; summing over $\Nf$ Dirac species and integrating
against the $\LambdaUV^{4}$ UV-momentum-volume factor reproduces our
$a_{0}(\Nf) \cdot \LambdaUV^{4} = (\Nf/4\pi^{2})\,\LambdaUV^{4}$.
This produces the emergent
cosmological-constant prediction
\begin{equation}
  \Lambdaemerg(\LambdaUV,\Nf) \;=\; a_{0}(\Nf)\cdot\LambdaUV^{4}
  \;=\; \frac{\Nf}{4\pi^{2}}\,\LambdaUV^{4}.
  \label{eq:Lemerg}
\end{equation}
The microscopic Newton constant under the ADW substrate-rescaling is
\begin{equation}
  \GN^{\mathrm{micro}}(\LambdaUV,\Nf,\aADW)
    \;=\; \aADW \cdot \GNsak(\LambdaUV,\Nf),
  \label{eq:GNmicro}
\end{equation}
where $\aADW$ is the dimensionless ADW coefficient
\cite{Vergeles2025}; at the Sakharov-Adler calibration $\aADW=1$
this matches the Phase~6a.1 emergent Newton constant
$\GNem(\LambdaUV,\Nf,1)=\GNsak$ \cite{LinearizedEFE2026}.

The match residual measures the algebraic discrepancy between the
ADW-rescaled microscopic Newton constant and the heat-kernel
baseline,
\begin{equation}
  \delta\GN(\LambdaUV,\Nf,\aADW) \;=\;
    (\aADW - 1)\cdot \GNsak(\LambdaUV,\Nf).
  \label{eq:dGN}
\end{equation}

\subsection{Observed cosmological constant}

The Planck-2018 measurement of the energy budget gives
\cite{Planck2018}
\begin{equation}
  \Lambdaobs \;\simeq\; (2.26\,\mathrm{meV})^{4}
    \;\simeq\; 2.6\times 10^{-47}\,\mathrm{GeV}^{4},
\end{equation}
encoded in our Lean module as the rational
$\Lambdaobs = 26/10^{48}\,\mathrm{GeV}^{4}$ and in the Python pipeline
as \texttt{LAMBDA\_OBSERVED\_GEV4} (see \texttt{constants.py}).

\section{Decision Gate E.4 quantitative theorem}

The substantive load-bearing claim is a Lean theorem:

\begin{theorem}[Decision Gate~E.4 quantitative anchor]
\label{thm:DG-E4}
At the natural microscopic point
$(\LambdaUV,\Nf) = (1.2\times 10^{19}\,\mathrm{GeV},\ 16)$, the
heat-kernel-induced emergent cosmological constant exceeds
$10^{100}\Lambdaobs$:
\begin{equation}
  \Lambdaemerg(\MPl,16) \;>\; 10^{100}\Lambdaobs.
\end{equation}
\end{theorem}

The proof uses the standard $\pi < 3.15$ bound (Mathlib
\texttt{Real.pi\_lt\_d2}) to establish $(4\pi)^{2} < 160$, hence
$a_{0}(16) > 4/10$; combined with $\MPl^{4} \simeq 2.07\times 10^{76}$
this yields $\Lambdaemerg > 8.3\times 10^{75}\,\mathrm{GeV}^{4}$ to
be compared with $10^{100}\Lambdaobs = 26\times 10^{52}\,\mathrm{GeV}^{4}$.
The Lean theorem \texttt{lambdaEmerg\-Microscopic\_at\_planck\_natural\_far\_exceeds\_observed}
is closed via this chain in
\texttt{lean/SKEFTHawking/MicroscopicCoefficientMatch.lean}.

\subsection{Numerical verdict at the natural-cutoff point}

Computed in the companion Python pipeline at full numerical
precision:
\begin{equation}
  \frac{\Lambdaemerg(\MPl,16)}{\Lambdaobs}
    \;\simeq\; 3.5\times 10^{122},
\end{equation}
i.e. the ratio is approximately $10^{122}$, recovering the standard
``120 orders of magnitude'' phrasing of the cosmological-constant
problem.

The Lean theorem proves $> 10^{100}$ (with margin); the actual
numerical value is $\simeq 10^{122}$ --- the Lean bound is the
quantitative \emph{floor} below which any heat-kernel-derivable ratio
in the SM-like band cannot fall.

\subsection{Verdict-band map}

We define three Decision-Gate-E.4 verdict bands:
\begin{itemize}
\item \texttt{cc\_resolved}:
  $|\log_{10}(\Lambdaemerg/\Lambdaobs)| < 1$
  (within an order of magnitude of $\Lambdaobs$).
\item \texttt{cc\_reproduced}: $\Lambdaemerg/\Lambdaobs > 10^{60}$
  (the classical CC problem in emergent form).
\item \texttt{cc\_intermediate}: neither resolved nor reproduced.
\end{itemize}

The natural high-energy theory band $\LambdaUV \geq 246\,\mathrm{GeV}$
(EW scale), $\Nf \geq 1$ sits entirely in \texttt{cc\_reproduced}
across the scan. The diagnostic resolution locus
$\LambdaUV \simeq 2.83\times 10^{-12}\,\mathrm{GeV} = 2.83\,\mathrm{meV}$
(four-meV scale, far sub-electroweak) is the unique unphysical band
where $\Lambdaemerg \simeq \Lambdaobs$ holds. No natural UV
completion produces a cutoff at the meV scale.

\begin{figure*}[t]
  \centering
  \includegraphics[width=\textwidth]{figures/fig_lambda_emerg_parameter_scan.png}
  \caption{Phase~6e Wave~5: $\Lambdaemerg$ parameter scan and Decision
  Gate~E.4 verdict bands. Left panel: log-log $\Lambdaemerg$ vs.
  $\LambdaUV$ at fixed $\Nf$ (curves for $\Nf = 1, 4, 16, 100$); the
  Planck-2018 reference $\Lambdaobs$ is drawn as a horizontal dotted
  anchor; the natural cutoff $\MPl$ is drawn as a dashed vertical
  anchor; the diagnostic resolution locus
  $\LambdaUV \simeq 2.83\times 10^{-12}\,\mathrm{GeV}$ is marked. Right
  panel: 2D verdict map of $\log_{10}(\Lambdaemerg/\Lambdaobs)$ over
  the $(\LambdaUV,\Nf)$ plane with $|\log_{10}\mathrm{ratio}|=1$
  (\texttt{cc\_resolved} boundary) and $\log_{10}\mathrm{ratio}=60$
  (\texttt{cc\_reproduced} boundary) contours overlaid. The natural
  high-energy theory band sits entirely in \texttt{cc\_reproduced}.}
  \label{fig:Le-scan}
\end{figure*}

\section{Microscopic-macroscopic match: $\delta\GN = 0$ biconditional}

Wave~5's secondary correctness-push (the Wave~5 expression of
Decision Gate~E.2 \cite{Roehm2026Wave1}) is the biconditional

\begin{theorem}[Match-residual biconditional]
\label{thm:match-bicond}
For positive $(\LambdaUV,\Nf)$, the match residual
$\delta\GN(\LambdaUV,\Nf,\aADW) = 0$ if and only if
$\aADW = 1$.
\end{theorem}

The forward direction uses $\GNsak > 0$ (Lean theorem
\texttt{G\_N\_from\_a2\_pos}, Wave~1) to flip the
$\mathrm{mul\_eq\_zero}$ split; the reverse is a definitional unfold.
The Lean theorem
\texttt{matchResidual\_eq\_zero\_iff\_alpha\_unity}
closes by these two directions in our module.

The substantive cross-bridge to Phase~6a.1 is the calibration witness
\texttt{gNMicroscopic\_at\_alpha\_one\_eq\_G\_N\_emerg}: at $\aADW=1$,
the microscopic Newton constant equals
$\GNem(\LambdaUV,\Nf,1)$. The Lean proof body invokes both
\texttt{LinearizedEFE.G\_N\_emerg\_at\_alpha\_one}
\cite{LinearizedEFE2026} and
\texttt{HeatKernelExpansion.G\_N\_from\_a2\_eq\_G\_N\_sakharov}
\cite{Roehm2026Wave1} by name --- a load-bearing P6 cross-module
bridge integrity per
\texttt{feedback\_python\_lean\_refs\_drift.md}.

\section{Stelle-basis higher-curvature aggregate}

The Wave~2 Stelle-basis triple $(\alpha,\beta,\gamma)$
\cite{HigherCurvature2026,Stelle1977}, parameterizing the order-$a_{4}$
density via
\[
  a_{4}\;=\;\alpha\, R^{2}+\beta\,C^{2}+\gamma\,\mathcal{G},
\]
satisfies the closed-form aggregate identity
\begin{equation}
  \alpha(\Nf)+\beta(\Nf)+\gamma(\Nf)
    \;=\; -\frac{7\,\Nf}{810}\cdot\frac{1}{(4\pi)^{2}},
  \label{eq:stelle-sum}
\end{equation}
established by the Lean theorem
\texttt{higherCurvature\_stelle\_sum\_eq}. Combining with
$\Nf > 0$ gives the strict-negativity aggregate
\texttt{higherCurvature\_stelle\_sum\_negative}
($\alpha+\beta+\gamma < 0$), which serves as the third conjunct of
the bundled match predicate.

\section{Bundled tracked-Prop match witness}

We define the bundled predicate
\begin{align}
&\mathtt{H\_MicroscopicCoefficientMatch}(\LambdaUV,\Nf,\aADW) :=\nonumber\\
&\quad \delta\GN(\LambdaUV,\Nf,\aADW) = 0\ \wedge\nonumber\\
&\quad 0 < \Lambdaemerg(\LambdaUV,\Nf)\ \wedge\nonumber\\
&\quad \alpha(\Nf)+\beta(\Nf)+\gamma(\Nf) < 0
\end{align}
encoding the Wave~5 microscopic-macroscopic match as a single
Prop consumable downstream. The Dirac witness theorem
\texttt{dirac\_H\_MicroscopicCoefficientMatch\_at\_alpha\_one} closes
it under positive $(\LambdaUV,\Nf)$ at $\aADW = 1$; the falsifier
\texttt{perturbed\_alpha\_not\_H\_MicroscopicCoefficientMatch} rules
out the bundle for any $\aADW \neq 1$.

\section{Discussion: shelter in emergent gravity?}

A common motivation for emergent-gravity programmes is the hope that
emergent formulations might \emph{naturally} produce a cosmologically
small $\Lambda$. Wave~5 closes that hope cleanly for the ADW
programme:
the heat-kernel $a_{0}$ contribution is dimension-correct
($\sim \LambdaUV^{4}$) by construction, and at any natural UV cutoff
the resulting $\Lambdaemerg$ is fine-tuning-distance away from
$\Lambdaobs$ in the same sense as in standard QFT. The
emergent-gravity reformulation does \emph{not} provide shelter from
the Weinberg cosmological-constant problem.

This is consistent with two related project results:
(i) Phase~5y closure of the dark-sector classification table found
that DE--mechanism candidates fail empirical falsifiers across the
q-theory and ADW landscapes; (ii) Phase~6c Wave~1 strong-CP
topological dark energy achieves $\Lambdaobs$ via a separate
mechanism (Zhitnitsky topological-charge density) at PDG
$\Lambda_{\mathrm{QCD}}$. Wave~5 then constrains the project's CC
mechanism to a \emph{single} channel: the heat-kernel
$a_{0}$ does not produce $\Lambdaobs$ naturally; the Phase~6c
topological-DE mechanism is the canonical CC source.

\section{Method --- Lean formalization}

The module
\texttt{lean/SKEFTHawking/MicroscopicCoefficientMatch.lean}
is verified by \texttt{lake build} (zero sorry, zero new
axioms). It contains \microscopicCoefficientMatchThms{} substantive
theorems (the count macro is auto-generated by
\texttt{scripts/update\_counts.py} from the canonical Lean source).
The companion Python pipeline at
\texttt{src/micro\_macro\_match/} ships
\microscopicCoefficientMatchTests{} pytest cases (all passing).
The build is reproducible from a clean baseline by
\texttt{rm -rf .lake/build \&\& lake build SKEFTHawking.ExtractDeps}.

The single substantive cross-bridge to Phase~6a.1 invokes
\texttt{LinearizedEFE.G\_N\_emerg\_at\_alpha\_one}; the
cross-bridge to Wave~1 invokes
\texttt{HeatKernelExpansion.G\_N\_from\_a2\_eq\_G\_N\_sakharov};
the cross-bridge to Wave~2 unfolds the Wave~2 definitions
\texttt{a4\_alpha}, \texttt{a4\_beta}, \texttt{a4\_gamma} (P6
drift-protection per
\texttt{feedback\_python\_lean\_refs\_drift.md}) and closes by
\texttt{ring} on the rational aggregate (theorem
\texttt{higherCurvature\_stelle\_sum\_eq}); the sign-aggregate
corollary \texttt{higherCurvature\_stelle\_sum\_negative}
discharges by \texttt{nlinarith} on \texttt{fourPiSqInv\_pos}.

\section{Conclusion}

Wave~5 closes Decision Gate~E.4 with the verdict
\texttt{cc\_reproduced} at the natural microscopic point: the
ADW heat-kernel programme reproduces the cosmological-constant
problem in emergent form, with no fine-tuning shelter under natural
$(\LambdaUV,\Nf)$. The match residual
$\delta\GN(\LambdaUV,\Nf,\aADW)$ vanishes iff the ADW substrate
parameter is at the Sakharov-Adler calibration $\aADW = 1$, providing
the Wave~5 expression of the Decision Gate~E.2 closure at the
microscopic-coefficient level. Phase~6e Wave~6 (Einstein-Cartan
torsion-sourced extension) consumes the bundled
\texttt{H\_MicroscopicCoefficientMatch} as the input matching
condition.

\begin{thebibliography}{99}

\bibitem{Sakharov1968}
A. D. Sakharov, ``Vacuum quantum fluctuations in curved space and the
theory of gravitation'', \emph{Sov. Phys. Dokl.} \textbf{12},
1040 (1968).

\bibitem{Adler1982}
S. L. Adler, ``Einstein gravity as a symmetry-breaking effect in
quantum field theory'', \emph{Rev. Mod. Phys.} \textbf{54}, 729
(1982).

\bibitem{Diakonov2011}
D. I. Diakonov, ``Towards lattice-regularized quantum gravity'',
\emph{arXiv:1109.0091} (2011).

\bibitem{VladimirovDiakonov2012}
A. A. Vladimirov and D. I. Diakonov, ``Phase transitions in spinor
quantum gravity on a lattice'', \emph{Phys. Rev. D} \textbf{86},
104019 (2012).

\bibitem{Vergeles2025}
S. N. Vergeles, ``Unitarity of 4D Lattice Theory of Gravity'',
\emph{Phys. Rev. D} \textbf{112}, 054509 (2025).

\bibitem{Vassilevich2003}
D. V. Vassilevich, ``Heat kernel expansion: user's manual'',
\emph{Phys. Rep.} \textbf{388}, 279 (2003).

\bibitem{Stelle1977}
K. S. Stelle, ``Renormalization of higher-derivative quantum
gravity'', \emph{Phys. Rev. D} \textbf{16}, 953 (1977).

\bibitem{Planck2018}
N. Aghanim \emph{et al.} (Planck Collaboration), ``Planck 2018
results. VI. Cosmological parameters'', \emph{A\&A} \textbf{641},
A6 (2020).

\bibitem{LinearizedEFE2026}
J. G. Roehm, ``Sakharov-Adler induced gravity from the ADW emergent
programme: a Lean-formalized linearized Einstein field equation'',
in preparation (2026).

\bibitem{Roehm2026Wave1}
J. G. Roehm, ``Formal heat-kernel calibration of induced Newton's
constant from the ADW Dirac substrate'', in preparation (2026).

\bibitem{HigherCurvature2026}
J. G. Roehm, ``Higher-curvature Stelle-basis structure from the
Christensen-Duff Dirac sector at order $a_{4}$'', in preparation
(2026).

\bibitem{Roehm2026Wave4}
J. G. Roehm, ``Variational nonlinear Einstein field equations from
the ADW emergent-gravity programme'', in preparation (2026).

\end{thebibliography}

\end{document}
